%\begin{theorem}[Main Runtime Bound] \label{th:expoHD-cliff}
%	Let $d < n(\tfrac14-\varepsilon)$ for fixed $\varepsilon>0$, and let the parameter $c$ of the IPH 
%	with linear mutation potential  $M(x)=\max\{1,\lfloor c\,\mathrm{HD}(x,\text{best})\rfloor\}$ 
%	for some 
%	\[
%	c < \frac{1+4\varepsilon}{3-4\varepsilon}.
%	\]
%	If the ageing threshold satisfies $\tau=\Omega(n^{1+\varepsilon'})$ for some 
%	$\varepsilon'>0$, then the $(1+1)$~Opt-IA$_{\textsc{FirstBest}}$ optimises \cliff{} in expected
%	\[
%	O\!\left(\frac{n^{5}\log(n/d)}{d^{3}} + n^{2}\log d\right)
%	\]
%	fitness evaluations.
%\end{theorem}

\rev{The following theorem is the main result of this section.}
\begin{theorem}\label{th:expoHD-cliff}
	The %{\oneoneOPTIA} using \IPHfcm{} 
	\textcolor{black}{\firstBest~with} \linHD\ (factor $c<\tfrac{1+4\varepsilon}{3-4\varepsilon}$) and ageing parameter $\tau=\Omega(n^{1+\epsilon'})$
	optimises \cliff{} with $d<n(\tfrac{1}{4}-\epsilon)$ in expected
	\[
	O\!\Big(\frac{n^{5}\log(n/d)}{d^{3}} \;+\; n^{2}\log d\Big)
	\]
	fitness function evaluations.
\end{theorem}


We now present the proof sketch of our main result on the  \cliff{} function, together with 
all supporting lemmata.  %\textcolor{green}{All proofs are deferred to...} %the appendix; the 
%statements here are mathematically identical to those in the full technical 
%development, but streamlined for readability.  
The argument proceeds in two 
major stages: (i) obtaining a canonical entry point onto the second slope 
immediately after an ageing event at the cliff-top; and (ii) showing that from 
this configuration the algorithm reaches the optimum without dangerous increases in mutation potential.

For stage (ii) we will use the generalization of Markov's Inequality by Ville~\cite{Doob1953}(Ch.~VII) which will allow us to show that the mutation potential stays bounded with constant probability.
\begin{theorem}[Ville's inequality for nonnegative supermartingales]\label{thm:ville}
	Let $(M_t)_{t\ge 0}$ be a nonnegative supermartingale adapted to $(\mathcal F_t)_{t\ge 0}$. 
	Then for every $a>0$,
	\[
	\Pr\!\big(\sup_{t\ge 0} M_t \ge a\big)\ \le\ \frac{\EE[M_0]}{a}.
	\]
	%\emph{(Ville, 1939; see also Doob, 1953, Ch.~VII.)}
\end{theorem}



To establish the result, we first analyse how the algorithm reaches a highly 
structured ``canonical'' state shortly after ageing triggers at the cliff-top.
At the cliff-top the algorithm cannot improve except by sampling the global 
optimum, and so ageing will eventually be triggered.  The key insight is that 
the iteration of ageing followed by two additional \textcolor{black}{fitness improvements} can place the 
algorithm into a
\textcolor{blue}{canonical
starting configuration on the second slope.
Throughout the \cliff{} analysis we index the successful fitness improvements on
the second slope by $r = r_0, r_0+1, \dots$, and write $\mathrm{HD}(r)$ for
the Hamming distance between the current solution and \emph{best} after $r$
such improvements.
With this notation, the canonical start corresponds to
$r_0 = 2$ (two improvements completed from the cliff bottom)
and $\mathrm{HD}(r_0) = 2$ (the current solution differs from the locked
\emph{best} in exactly two positions).}  The following lemma formalises this.

\begin{lemma}%[Canonical Start After Cliff-Top Ageing]
	\label{lem:can-start-omega-d3}
	Suppose ageing is triggered at the cliff-top ($d$~zero-bits).  
	The \rev{\firstBest}~with \linHD\ with $c\in(0,1)$ reaches within at most three steps the 
	configuration
	\[
	r_0=2, \qquad \mathrm{HD}(r_0)=2,
	\]
	with probability at least 
	\[
	\frac{3}{256}\left(\frac{d}{n}\right)^3.
	\]
\end{lemma}

\begin{proof}
	We exhibit a consecutive three-step event which happens with probability $\Omega((d/n)^3)$ and yields the stated starting 
	position on the second slope.
	
	 In the ageing iteration we require that the first two flips are $1\!\to 0$ then $0\!\to 1$ without replacement. This has probability $\frac{n-d}{n}\cdot\frac{d}{n-1}$ and produces an \emph{equal-fitness} offspring accepted by FCM (first return to $\ge$ parent). The offspring inherits the parent's age (no improvement), so both parent and offspring face the ageing coin and die independently with probability $1/2$. Requiring "parent dies, offspring survives'' contributes a factor $1/4$, and by the update rule the reference locks to the last cliff-top (the parent that died). The surviving equal-fitness offspring differs from this locked reference in exactly two bits, hence $\HD=2$ and, since $c<1$, the next iteration has $M=1$.
		
 With $M=1$, exactly one flip is performed and FCM does not accept worse-than-parent, so the algorithm returns the single-flip string after the potential exhaustion. A $0\!\to 1$ flip occurs with probability at least $d/n$ and returns the \emph{cliff-bottom} ($n-d+1$ ones), at Hamming distance $1$ from the locked reference. Ageing does not reset (no improvement), so both the current parent (equal-fitness cliff-top) and the returned offspring face the coin; requiring "parent dies, offspring survives'' contributes another factor $1/4$ and preserves the locked reference at the last cliff-top.
	
 In the following iteration with probability at least $(d-1)/n$ the first bit flip is a $0\!\to 1$ flip, and the offspring is \emph{accepted} by FCM (strict improvement from the cliff-bottom), producing a point with $n-d+2$ ones. Relative to the locked cliff-top reference, the Hamming distance is now $2$, and exactly two successful phases have completed: $r_0=2$, $\HD(r_0)=2$.
	
	Multiplying the lower bounds yields
	\[
	\Big(\tfrac{n-d}{n}\cdot\tfrac{d}{n-1}\cdot\tfrac14\Big)\cdot
	\Big(\tfrac{d}{n}\cdot\tfrac14\Big)\cdot
	\Big(\tfrac{d-1}{n}\Big).
	\]
	For $d\le(\tfrac14-\varepsilon)n$ we have $(n-d)/n\ge 3/4$, and for $d\ge 2$ we have $\tfrac{d}{n-1}\ge\tfrac{d}{n}$ and $\tfrac{d-1}{n}\ge \tfrac12\cdot\tfrac{d}{n}$. Hence the product is at least
	\[
	\frac{3}{256}\,\Big(\frac{d}{n}\Big)^{\!3},
	\]
	which proves the claim.
\end{proof}

The canonical condition $(r_0,\mathrm{HD}(r_0)) =
(2,2)$ means that we begin the drift analysis two successful improvements into
the second slope, at a point differing from the reference in exactly two
bit positions.

Repeated ascents to the cliff-top therefore form a geometric waiting-time 
process.   \rev{It takes at most $\sum_{i=d}^{n}n/i=O(n\log(n/d))$ for RLS to climb the first slope and the mutation potential is at most linear.
	 %Thus, each ascent takes $O(n^2\log(n/d))$ evaluations in expectation.
	 } Since each ascent takes $\rev{O}(n^2\log(n/d))$, the 
total time to obtain the canonical start is \new{provided by the following corollary}.

\begin{corollary}
	%[Expected Time to Canonical Start]
	The expected evaluation count until the first occurrence of 
	$(r_0,\mathrm{HD}(r_0)) = (2,2)$ is 
	\[
	O\!\left(\frac{n^{5}\log(n/d)}{d^{3}}\right).
	\]
\end{corollary}


From the canonical starting point, the key difficulty is \new{showing that} the mutation
potential \new{does not increase} too much.  With only few zero-bits remaining, successful
steps on the second slope require flipping one of these rare zeros, whereas each
$1\!\to0$ flip increases the Hamming distance and thus the mutation potential.
If this potential grows too large, hypermutations become long enough that, under
FCM$(\ge)$, they are increasingly likely to wander back to the cliff-top; once
there, the stopping rule causes the process to lock onto the plateau and erase
all progress.  Hence the analysis must show that with constant probability the
mutation potential stays small throughout the entire second-slope climb.  This
requires  controlling tightly how many detrimental flips a hypermutation can
accumulate before recovering to parent-level fitness.  The next lemma provides
the crucial starting point: an exponential tail bound for this damaging prefix.

\begin{lemma}%[Exponential Tail for Catch-Up Length]
	 \label{lem:xi-tail}
	\textcolor{blue}{Let the current solution on the second slope have $z\le d$ zero-bits and
	define $p=z/n$.
	Consider a single application of the hypermutation operator; let $\xi_z$ be
	the number of $1\to0$ flips made \emph{within that one hypermutation} before
	the fitness first returns to $\ge$ the parent's fitness (i.e., before the
	first catch-up).}
	Then for all $k\ge1$,
	\[
	\Pr(\xi_z\ge k)\le (2p)^k.
	\]
	Consequently, the probability that the Hamming distance increases by at least 
	$2k$ in a single hypermutation is also at most $(2p)^k$.
\end{lemma}

\begin{proof}
	If the first flip is a $0\!\to 1$ move, the offspring is strictly better and the hypermutation stops immediately, yielding $\xi_z=0$. Hence the event $\{\xi_z\ge 1\}$ implies that the first flip was a $1\!\to 0$ move. From that point on the process carries a \emph{deficit} of one-bits $D_t:=\#(1\!\to 0)-\#(0\!\to 1)$ that starts at $D_1=1$ and decreases by $1$ on each $0\!\to 1$ flip and increases by $1$ on each $1\!\to 0$ flip. The hypermutation stops at the \emph{first} time $t$ with $D_t=0$ (a catch-up). Thus the event $\{\xi_z\ge k\}$ is exactly the event that the first return to $0$ happens no earlier than time $2k$.
	
	We bound this probability from above by a simple ``coarse blocking'' argument that uses only that $z$ is small and that flips occur without replacement.
	
	\smallskip
	Condition on any history of length $t\le 2k-1$ that keeps $D_t\ge 1$ (so no catch-up yet) and in which at most $t$ distinct positions have been flipped. Among the $n-t$ remaining positions there are at most $z$ zeros. Hence the conditional probability that the \emph{next} flip is a $0\!\to 1$ move is at most
	\[
	\frac{z}{\,n-t\,}\ \le\ \frac{z}{\,n-(2k-1)\,}\ \le\ \frac{z}{(1/2+2\varepsilon)n}\ \le\ 2p,
	\]
	because $2k\le 2d<(\tfrac12-2\varepsilon)n$. In other words, throughout any hypermutation that lasts up to $2k$ flips, a zero-flip is never more likely than $2p$.
	
	\smallskip
	If the first catch-up occurs at some time $\ge 2k$, then along the way the trajectory must realize \emph{at least $k$} distinct $0\!\to 1$ flips while $D>0$ (each such flip reduces the deficit by $1$ but, by assumption, does not yet hit $0$ before the $k$-th such zero\new{-flip}). Expose these $k$ zero-flips in order and ignore everything else (the interspersed $1\!\to 0$ flips only make catch-up \emph{harder} and so can be dropped for an upper bound). By Step~1, each of these $k$ zero\new{-flips} occurs with conditional probability at most $2p$, uniformly over the entire run while $D>0$. Therefore, by the chain rule,
	\[
	\Pr(\text{first catch-up occurs at time }\ge 2k)\ \le\ (2p)^k.
	\]
	
	
	Since the left-hand side is exactly $\Pr(\xi_z\ge k)$, the claim follows.
\end{proof}


A whole \emph{phase} consists of all attempts at the same zero-bit count until the 
first success.  
Using the tail bound above, we obtain uniform exponential moment bounds for 
entire phases.

\begin{lemma}%[Phase Moment-Generating Bound] 
	\label{lem:phase-mgf}
	Consider a phase at zero-bit count $z$ and let $Y_z$ denote the total 
	Hamming-distance increase from all failed attempts in the phase.  
	Let $\rho = 2d/n$.  
	For all $0<\lambda<\tfrac12\ln\frac{1}{2\rho}$,
	\[
	\mathbb{E}[e^{\lambda Y_z}] 
	\;\le\; 
	\textcolor{blue}{\frac{1-\rho}{1-\rho e^{2\lambda}}
	\;\eqqcolon\; K(\lambda).}
	\]
\end{lemma}

\begin{proof}
	For a single \emph{failed} attempt starting the hypermutation with $z$ zeroes, let $\xi$ count the number of $1\!\to 0$ flips while below the parent's fitness before the first catch-up.
	By Lemma~\ref{lem:xi-tail} and the change in Hamming distance $\Delta_{\mathrm{fail}}\le 2\xi$ with $\xi\ge 1$,
	\begin{align*}
		\EE\!\left[e^{\lambda \Delta_{\mathrm{fail}}}\right]
		&\le \sum_{k\ge 1}^{\infty} e^{2\lambda k}\,\Pr(\xi=k)
		\\ &\le\ \sum_{k\ge 1}^{\infty} \big(2p\,e^{2\lambda}\big)^{k}
		\\ &\textcolor{black}{=}\ \frac{2p\,e^{2\lambda}}{\,1-2p\,e^{2\lambda}\,}:=:\ t.
	\end{align*}
	Let $F$ be the number of failed attempts before the first success in the phase. %with the convention
	%$\Pr(F=f)=(1-p_z)^f p_z$ for $f\in\{0,1,\ldots\}$ and $p_z\ge p=z/n$.
	Since $Y_z$ is the sum over failed attempts only (the success contributes $0$), conditioning stepwise on the phase history gives
	\begin{align*}
		\EE\!\big[e^{\lambda Y_z}\big]
		&=\EE\!\Big[\ \prod_{i=1}^{F} e^{\lambda \Delta_{\mathrm{fail}}^{(i)}}\ \Big]
		\ \le\ \EE\!\big[t^{F}\big].
	\end{align*}
	Using the geometric MGF for $t\in[0,1)$,
	\begin{align*}
		\EE\!\big[t^{F}\big]
		&=
		 \sum_{f\ge 0} (1-p)^f p\, t^f
		\ =\ \frac{p}{\,1-(1-p)t\,}.
	\end{align*}
	Substituting $t=\dfrac{2p\,e^{2\lambda}}{1-2p\,e^{2\lambda}}$ and simplifying,
	\begin{align*}
		\EE\!\big[e^{\lambda Y_z}\big]
		&\le \frac{p}{\,1-(1-p)\,\dfrac{2p\,e^{2\lambda}}{1-2p\,e^{2\lambda}}\,}
		\\ &=\ \frac{p\,(1-2p\,e^{2\lambda})}{\,1-2p\,e^{2\lambda}(2-p)\,}
		\ :=:\ K_z(\lambda).
	\end{align*}
	Since $p\le \rho/2$ and the map $p\mapsto K_z(\lambda)$ is increasing on $(0,\tfrac12)$ for fixed $\lambda\in(0,\tfrac12\ln(1/\rho))$, we obtain the uniform bound
	\begin{align*}
		\EE\!\big[e^{\lambda Y_z}\big]
		\ \le\ K_z(\lambda)\Big|_{p=\rho/2}=\frac{\tfrac{\rho}{2}(1-\rho\,e^{2\lambda})}{1-\rho e^{2\lambda}\,(2-\tfrac{\rho}{2})}
	\end{align*}
	Now, our claimed upper bound reduces to showing that a concave quadratic function of $u:=\rho e^{2\lambda} \in \left[0,\frac{1}{2}\right]$ is non-negative.
	\begin{align*}
		&\frac{\tfrac{\rho}{2}(1-\rho\,e^{2\lambda})}{1-\rho e^{2\lambda}\,(2-\tfrac{\rho}{2})}
		\  \le\ \frac{1-\rho}{\,1-\rho\,e^{2\lambda}\,}
		\ =:\ K(\lambda),
		\\&\Longleftrightarrow\quad
		\frac{\rho}{2}(1-u)^2 \le (1-\rho)\Bigl(1-\bigl(2-\tfrac{\rho}{2}\bigr)u\Bigr),\ \ u:=\rho\,e^{2\lambda},
		\\
		D(u)
		&:=(1-\rho)\Bigl(1-\bigl(2-\tfrac{\rho}{2}\bigr)u\Bigr)-\frac{\rho}{2}(1-u)^2
		\\&= \Bigl(1-\tfrac{3}{2}\rho\Bigr)+\Bigl(-2+\tfrac{7}{2}\rho-\tfrac{1}{2}\rho^2\Bigr)u-\tfrac{\rho}{2}u^2,
		\\
		D(0)&=1-\tfrac{3}{2}\rho\ \ge\ \tfrac14,\qquad
		D\!\Big(\tfrac12\Big)=\tfrac{\rho}{8}\bigl(1-2\rho\bigr)\ \ge\ 0
		\quad\Rightarrow\quad \\
		D(u)&\ge 0\ \text{ for }u\in[0,\tfrac12],
	\end{align*}
	
	
\end{proof}


To track the safety of the mutation potential across phases, we introduce a 
linear \emph{buffer} \new{$B_r = \theta r - s_0 - E_r$} capturing the invariant $c\,\mathrm{HD}(r)\le r$.

\begin{lemma}%[Buffer Process and Safety Criterion] 
	\label{lem:buffer}
	Let successful phases be indexed by $r=r_0,r_0+1,\dots$ and define
	\[
	s_0=\mathrm{HD}(r_0)-r_0, \qquad 
	\theta=\frac{1}{c}-1, \qquad
	E_r=\sum_{j=r_0+1}^{r} Y_{z_j},
	\]
	and $B_r = \theta r - s_0 - E_r$.
	Then for all $r\ge r_0$:
	\begin{enumerate}
		\item 
		$B_r\ge0$ if and only if $c\cdot\mathrm{HD}(r)\le r$;
		\item 
		$B_{r+1}-B_r=\theta-Y_{z_{r+1}}$;
		\item 
		$B_{r_0}=\theta r_0 - s_0$.
	\end{enumerate}
\end{lemma}

\begin{proof}
	For any $r\ge r_0$, the Hamming distance admits the exact decomposition
	\begin{align*}
		\HD(r)\ &=\ \underbrace{\HD(r_0)}_{=\,r_0+s_0}\ +\ \underbrace{(r-r_0)}_{\text{one per success}}\ +\ \underbrace{E_r}_{\text{ extra HD in failed attempts}}
		\ \\&=\ r+s_0+E_r,		
	\end{align*}
	
	since each strict improvement increases the HD by exactly $1$ and $E_r=\sum_{j=r_0+1}^r Y_{z_j}$ collects the extra HD from failed attempts (the successful attempt contributes $0$ because $\Delta_z^{\mathrm{succ}}=1$).
	
	\smallskip
	\emph{(i) Safety invariant.}
	Using $\theta=\tfrac{1}{c}-1$ and $B_r:=\theta r - s_0 - E_r$,
	\begin{align*}
		B_r\ge 0
		\ &\Longleftrightarrow\
		\theta r \ge s_0+E_r
		\ \Longleftrightarrow\
		c\theta r \ge c(s_0+E_r)
		\\ &\Longleftrightarrow\
		(1-c)r \ge c(s_0+E_r).	
	\end{align*}
	
	
	But $\HD(r)=r+s_0+E_r$, so the last inequality is equivalent to
	\[
	r \ \ge\ c\,(r+s_0+E_r)\ =\ c\,\HD(r),
	\]
	i.e., $c\,\HD(r)\le r$. Thus $c\,\HD(r)\le r \iff B_r\ge 0$.
	
	\smallskip
	\emph{(ii) One-phase update.}
	By definition,
	\begin{align*}
		B_{r+1}-B_r
		&= \bigl(\theta(r+1)-s_0-E_{r+1}\bigr) - \bigl(\theta r - s_0 - E_r\bigr)
		\\&= \theta - (E_{r+1}-E_r)
		= \theta - Y_{z_{r+1}}.
	\end{align*}
	
	
	
	\smallskip
	(iii) Trivial from the definition at $r=r_0$.
\end{proof}

Combining \rev{the MGF of the buffer growth} with the buffer update yields an exponential drift 
envelope suitable for applying \textcolor{black}{Ville's inequality for non-negative supermartingales (Theorem~\ref{thm:ville}).}

\begin{lemma}
	%[Exponential Drift Envelope] 
	\label{lem:drift-envelope-global}
	Let $\rho=2d/n$.  
	For any $0<\lambda<\tfrac12\ln\tfrac{1}{2\rho}$ define
	\[
	\textcolor{blue}{\frac{1-\rho}{1-\rho e^{2\lambda}} \;\eqqcolon\; K(\lambda),}
	\qquad 
	\alpha=e^{-\lambda\theta}K(\lambda).
	\]
	Then for all $t\ge0$,
	\[
	\mathbb{E}[e^{-\lambda(B_{t+1}-B_t)}\mid\mathcal{F}_t]\le \alpha.
	\]
\end{lemma}

\begin{proof}
	By Lemma~\ref{lem:buffer}, $B_{t+1}-B_t=\theta-Y_{z_{t+1}}$, hence
	\[
	e^{-\lambda(B_{t+1}-B_t)}\ =\ e^{-\lambda\theta}\,e^{\lambda Y_{z_{t+1}}}.
	\]
	By Lemma~\ref{lem:phase-mgf}, for $0<\lambda<\tfrac12\ln\!\tfrac{1}{2\rho}$ we have %the uniform envelope
	$\EE[e^{\lambda Y_{z_{t+1}}}\mid\mathcal F_t]\le K(\lambda)$ for all phases (all $z_{t+1}\le d$). Therefore
	\begin{align*}
		\EE\!\big[e^{-\lambda(B_{t+1}-B_t)}\mid \mathcal F_t\big]
		\ &=\ e^{-\lambda\theta}\,\EE\!\big[e^{\lambda Y_{z_{t+1}}}\mid\mathcal F_t\big]
		\ \\ &\le\ e^{-\lambda\theta}\,K(\lambda)
		\ =\ \alpha	
	\end{align*}
\end{proof}

\new{We now apply Ville's inequality to turn such a drift envelope into a bound on barrier crossing.}


\begin{corollary}
	%[Buffer Barrier Bound] 
	\label{cor:ville-exponential}
	
	If $B_0=b>0$ and there exist $\lambda>0$ and $\alpha<1$ with
	\[
	\mathbb{E}[e^{-\lambda(B_{t+1}-B_t)}\mid\mathcal{F}_t]\le\alpha,
	\]
	then
	\[
	\Pr\!\left(\inf_{t\ge0} B_t \le 0\right) \le e^{-\lambda b}.
	\]
\end{corollary}

\begin{proof}
	\textcolor{blue}{Ville's inequality (Theorem~\ref{thm:ville}) applies to nonneg\-ative
	supermartingales, but the buffer $B_t$ is not one (it can drift and take
	arbitrary values).  We therefore work instead with the transformed process
	$M_t := \alpha^{-t}e^{-\lambda B_t}$, which \emph{is} a nonneg\-ative
	supermartingale: the assumption $\EE[e^{-\lambda(B_{t+1}-B_t)}\mid\mathcal F_t]
	\le \alpha$ gives $\EE[M_{t+1}\mid\mathcal F_t]\le M_t$.
	Applying Theorem~\ref{thm:ville} to $M_t$ with $a=1$ yields
	$\Pr(\sup_{t\ge 0} M_t\ge 1)\le \EE[M_0]=e^{-\lambda b}$.
	Finally, $B_t\le 0$ implies $e^{-\lambda B_t}\ge 1$, hence $M_t\ge 1$,
	giving the claim.}
	
\end{proof}


To guarantee the drift negativity required above, we identify the admissible 
range of the mutation-potential factor~$c$.

\begin{lemma}%[Per-Instance $c$-Threshold]
	 \label{lem:c-threshold-instance}
	
	Let $\rho=2d/n$ and $\theta=\frac{1}{c}-1$.  
	Define
	\[
	c_\star(n,d)=\frac{1-\rho}{1+\rho}.
	\]
	If $c<c_\star(n,d)$, then there exists $\lambda>0$ such that
	\[
	e^{-\lambda\theta}K(\lambda) < 1,
	\]
	and hence the buffer satisfies the Ville barrier bound.
\end{lemma}

\begin{proof}
	Let $g(\lambda):=\ln\alpha(\lambda)=-\lambda\theta+\ln\!\big(\tfrac{1-\rho}{1-\rho e^{2\lambda}}\big)$. 
	Then $g(0)=0$ and $g'(0)=-\theta+\tfrac{2\rho}{1-\rho}$. 
	If $\theta>\tfrac{2\rho}{1-\rho}$ (equivalently $c<\tfrac{1-\rho}{1+\rho}$), continuity gives $g(\lambda)<0$ for all sufficiently small $\lambda>0$, i.e., $\alpha(\lambda)<1$ in the admissible domain.
\end{proof}

A uniform instantiation over all permitted~$d$ shows that the feasible~$c$ 
range has a nontrivial lower bound.

\begin{corollary}%[Uniform Lower Bound for $c$]
	If $d\le (\tfrac14-\varepsilon)n$, then
	\[
	c_\star(n,d) \;\ge\; \frac{1+4\varepsilon}{3-4\varepsilon} \;\ge\; \tfrac13.
	\]
\end{corollary}
\label{cor:c-floor}



Finally, combining the threshold on $c$ with the initial buffer value at 
$(r_0,\mathrm{HD}(r_0))=(2,2)$ yields a uniform constant bound on the 
probability of avoiding jumping back.

\begin{lemma}
	%[Constant Safety Probability] 
	\label{lem:param-constant-c-general}
	If $c<c_\star(n,d)$ and 
	$b=B_{r_0}=\frac{r_0}{c}-\mathrm{HD}(r_0)>0$, 
	then 
	\[
	\Pr\!\left(\inf B_t \le 0\right) \le e^{-\lambda b}
	\]
	for some $\lambda>0$ depending only on $c$ and $d/n$.  
	For fixed~$\varepsilon$, this yields a uniform constant $<1$ independent of 
	$n,d$.
\end{lemma}

\begin{proof}
	By Lemma~\ref{lem:c-threshold-instance} (with $\rho=2d/n$), there exists
	$\lambda\in\big(0,\tfrac12\ln\tfrac{1}{2\rho}\big)=\big(0,\tfrac12\ln\tfrac{n}{4d}\big)$
	such that $e^{-\lambda\theta}K(\lambda)<1$. Applying Corollary~\ref{cor:ville-exponential}
	to $B_t$ started at $B_{r_0}=b>0$ yields
	$\Pr(\inf_{t\ge 0} B_t\le 0)\le e^{-\lambda b}$.
\end{proof}

Once the mutation potential remains safe throughout, progress on the second 
slope is monotone and its runtime is easily bounded.



\begin{lemma}%[Second-Slope Climb Time] 
	\label{prop:climb-time}
	
	If $B_t>0$ for all $t$, then the second-slope ascent completes in expected
	\[
	\Theta(n\log d)
	\]
	iterations.  
	If $\tau=\Omega(n^{1+\varepsilon})$, then $\Theta(n\log d)=o(\tau)$.
\end{lemma}

\begin{proof}
	When the zero-bit count is $z$, the success probability per attempt is at least $z/n$, so the expected number of attempts for $z\!\to\! z-1$ is at most $n/z$.
	Summing over $z=1,\dots,d$ gives $\sum_{z=1}^{d}n/z=\Theta(n\log d)$.
	Since $d\le(\tfrac14-\varepsilon)n$, we have $\log d=O(\log n)$, hence $n\log d=o(n^{1+\varepsilon})$ and $\EE[T_{\mathrm{climb}}]=o(\tau)$.
\end{proof}


The expected time to reach the canonical start is 
$O(n^{5}\log(n/d)/d^{3})$.  
Once there, with constant probability the buffer remains positive for all 
phases; conditioned on this event, the second-slope ascent costs 
$O(n^{2}\log d)$ evaluations.  
Since the success probability is constant, this remains the unconditional 
expectation.  
Summing the two contributions establishes the main runtime bound.

\begin{proof}[Proof [of Theorem~\ref{th:expoHD-cliff}]]
	
	From a uniformly random start, the expected iterations to reach the cliff-top ($n-d$ ones) are
	$\sum_{i=d}^{n}\frac{n}{i}=\Theta\!\big(n\log\tfrac{n}{d}\big)$, which at $\Theta(n)$ evaluations/iteration costs
	$\Theta\!\big(n^{2}\log\tfrac{n}{d}\big)$ evaluations per ascent to the cliff-top. At the cliff-top, the only fitter point is $1^n$, so ageing triggers unless the optimum is found within $\tau$ iterations.
	
	By Lemma~\ref{lem:can-start-omega-d3}, within at most three iterations after an ageing kick the process reaches the canonical second-slope start
	\[
	r_0=2,\qquad \HD(r_0)=2
	\]
	with probability
	\[
	p_{\mathrm{can}}(n,d)\ \ge\ \frac{3}{256}\Big(\frac{d}{n}\Big)^{\!3}.
	\]
	Hence the number of ageing epochs needed to obtain this start is geometric with mean $1/p_{\mathrm{can}}$, and the expected evaluation cost up to the canonical start (including restarts) is
	\[
	\Theta\!\Big(\frac{1}{p_{\mathrm{can}}(n,d)}\cdot n^{2}\log\tfrac{n}{d}\Big)
	\ =\ O\!\Big(\frac{n^{5}\log(n/d)}{d^{3}}\Big).
	\]
	
	\smallskip
	Under the reference-update rule the start on the second slope is $r_0=2$, $\HD(r_0)=2$, so $s_0=0$ and with $\theta=\tfrac1c-1$ define
	\[
	B_r\ :=\ \theta r - \sum_{j=r_0+1}^{r} Y_{z_j}\qquad(r\ge r_0),
	\]
	so $B_{r_0}=2\theta$ and $B_{t+1}-B_t=\theta-Y_{z_{t+1}}$ by Lemma~\ref{lem:buffer}.
	By the refined phase MGF (Lemma~\ref{lem:phase-mgf}), for every $0<\lambda<\tfrac12\ln\!\tfrac{1}{2\rho}$ with $\rho:=2d/n$,
	\[
	\textcolor{blue}{\EE\!\big[e^{\lambda Y_z}\big]\ \le\ \frac{1-\rho}{\,1-\rho e^{2\lambda}\,} \;\eqqcolon\; K(\lambda).}
	\]
	Therefore, by Lemma~\ref{lem:drift-envelope-global},
	\[
	\EE\!\big[e^{-\lambda(B_{t+1}-B_t)}\mid\mathcal F_t\big]\ \le\ e^{-\lambda\theta}\,K(\lambda)\qquad\text{for all }t\ge 0.
	\]
	
	Since $d\le(\tfrac14-\varepsilon)n$ and the theorem assumes
	\[
	c\ <\ \frac{1+4\varepsilon}{3-4\varepsilon}\ \le\ \frac{1-2d/n}{1+2d/n}\ =:\ c_\star(n,d),
	\]
	Lemma~\ref{lem:c-threshold-instance} guarantees an admissible $\lambda\in\big(0,\tfrac12\ln\tfrac{n}{4d}\big)$ with $e^{-\lambda\theta}K(\lambda)<1$.
	With $B_{r_0}=2\theta$, Corollary~\ref{cor:ville-exponential} yields
	\[
	\Pr\!\Big(\inf_{t\ge 0} B_t\le 0\Big)\ \le\ e^{-\lambda\,2\theta}.
	\]
	Moreover, since $\tfrac12\ln\!\tfrac{n}{4d}\ge \tfrac12\ln\!\tfrac{1}{1-4\varepsilon}$ for all admissible $d$, we may fix a \emph{uniform} $\lambda_\varepsilon\in\big(0,\tfrac12\ln\tfrac{1}{1-4\varepsilon}\big)$ and conclude
	\[
	\Pr\!\Big(\inf_{t\ge 0} B_t\le 0\Big)\ \le\ e^{-\lambda_\varepsilon\,2\theta}.
	\]
	Because $\lambda_\varepsilon$ depends only on $\varepsilon$ and $\theta=\tfrac1c-1$ depends only on $c$, this right-hand side is a constant in $(0,1)$, independent of $n$ and $d$.
	
	\smallskip
	On the event $\{\inf_{t\ge 0} B_t>0\}$ (no jump-back), the second-slope ascent is monotone and completes in expected
	\[
	\sum_{z=1}^{d}\frac{n}{z}\ =\ \Theta(n\log d)
	\]
	iterations, which is $o(\tau)$ since $\tau=\Omega(n^{1+\epsilon'})$. \textcolor{blue}{Since the mutation potential satisfies $M \le c\,n$ (where $c$ is the constant
	from the \linHD\ potential), each iteration evaluates at most $c\,n =\Theta(n)$
	individuals; one second-slope attempt therefore costs $\Theta(n^{2}\log d)$
	evaluations,} and the constant no-jump-back probability implies a constant expected number of such attempts.
	
	\smallskip

	Summing the (restarted) approach to the canonical start and the second-slope cost gives
	\[
	\EE[\text{evaluations}]
	\ =\ O\!\Big(\frac{n^{5}\log(n/d)}{d^{3}}\Big)\ +\ O\!\big(n^{2}\log d\big),
	\]
	as claimed.
\end{proof}

\textcolor{black}{The analysis for \lastBest~differs only in Lemma~\ref{lem:can-start-omega-d3} and actually yields a $d/n$ factor better probability bound since the mutation potential is $1$ at the cliff-top and thus, creating a solution at the bottom of the cliff is immediately possible when any cliff top individual dies, whereas for \firstBest~to work we have to first accept a neutral mutation to get the mutation potential to $1$, and then mutate into a cliff bottom solution in the next iteration. Consequently the runtime for \lastBest~is  
$O\!\Big(\frac{n^{4}\log(n/d)}{d^{2}} \;+\; n^{2}\log d\Big)$.}

\textcolor{black}{
\begin{corollary}
	The %{\oneoneOPTIA} using \IPHfcm{} 
	\textcolor{black}{\lastBest~with} \linHD\ \new{(factor $c<\tfrac{1+4\varepsilon}{3-4\varepsilon}$)} and ageing parameter $\tau=\Omega(n^{1+\epsilon'})$
	optimises \cliff{} with $d<n(\tfrac{1}{4}-\epsilon)$ in expected
	\[
	O\!\Big(\frac{\rev{n^{4}}\log(n/d)}{d^{2}} \;+\; n^{2}\log d\Big)
	\]
	fitness function evaluations.
\end{corollary}
}
The following section will highlight the advantages of using \textsc{First-Best} over using \textsc{Last-Best}.
